\documentclass[12pt]{article}
\usepackage{amsmath}
\usepackage{geometry}
\usepackage{setspace}
\geometry{margin=1in}
\setlength{\parindent}{0pt}
\setlength{\parskip}{1em}
\usepackage{hyperref}

\begin{document}

I have chosen to analyze two songs: Coldplay's \textit{``Viva La Vida"}[Ch15] and Radiohead's \textit{``Airbag"}[Ch14]. Both tracks exemplify their respective genres through stylistic choices, innovative instrumentation, and form.

\section*{Analysis of ``Viva La Vida" by Coldplay}
Coldplay's \textit{``Viva La Vida"} is representative of mainstream rock through its stylistic characteristics, instrumentation, and innovative use of form, reflecting the evolution of Britpop into a more universal, experimental sound. The piece opens with a ``driving string part that repeats a 4-chord sequence" (Covach 547), showcasing a melodic simplicity often seen in rock genres while employing orchestral strings, which add depth and grandeur. This creates a dynamic, cinematic feel that transcends typical rock instrumentation.

The song’s form, described as ``Compound AABA," highlights its mainstream accessibility. The repeating 8-measure string figure in the introduction and throughout most sections ``creates a sense of arc" while maintaining cohesion (Covach 547). This approach mirrors the genre's balance between innovation and listener familiarity. Instrumentation such as ``strings, kick drum, timpani, cymbals, synthesized piano sounds, bass, piano, lead and background vocals" (Covach 547) reflects the band’s shift towards a richer, more diverse sonic palette, blending classical and modern instruments.

Furthermore, production techniques emphasize dynamics and texture. The ``strong decrescendo and simplified texture" at the song's conclusion subvert expectations of a climactic ending, offering a reflective resolution (Covach 547). This nuanced approach aligns with mainstream rock's trend toward emotional depth and artistic exploration, moving beyond formulaic rock structures.

\section*{Analysis of ``Airbag" by Radiohead}
Radiohead's \textit{``Airbag"} exemplifies alternative rock through its intricate structure, innovative instrumentation, and seamless blend of traditional rock elements with experimental electronic sounds. The opening features ``a driving riff and a tambourine" with instruments entering sequentially, establishing a textured foundation (Covach 506). The use of ``ethereal instrumental melody" combined with an "active bass" during the verse reinforces the atmospheric qualities that define the track.

The song’s form is described as a ``simple verse with refrain," which, despite its straightforward label, allows for significant complexity. Each verse mirrors the structure of a 12-bar blues, with ``a statement, a restatement, and a contrasting statement" (Covach 506). The composed instrumental melody functions as ``an introduction, an interlude between verses, and a conclusive section," underscoring the song’s experimental nature and providing continuity throughout.

Instrumentation plays a key role in defining the song's identity within the genre. The use of ``drums, bass, electric guitars, various keyboards and synths, lead and backing vocals" (Covach 506) bridges the gap between traditional rock and electronic experimentation. The reliance on ``what seems to be a looped drum performance" integrates electronic music techniques into a rock context, creating a hypnotic rhythm that anchors the track amidst its complex arrangement.

The experimental nature of the song is further evident in the ``high-pitched ringing and other sound effects" that permeate the postlude, as well as the "experimental sounds" accompanying the extension section (Covach 506). These elements distinguish the song from conventional rock tracks and highlight Radiohead's progressive ethos, blending avant-garde noise with accessible melodic elements.

\begin{center}
\textbf{Works Cited}
\end{center}

Covach, John. \textit{What’s That Sound? An Introduction to Rock and Its History}. 5th ed., McGraw-Hill.
\end{document}
